\section{Initial Scope}

The first implementation should be intentionally small. The objective is not completeness. The objective is validating the abstraction.

\migaki{} v0 should include:

\begin{itemize}[leftmargin=1.5em]
  \item TypeScript runtime and fetch-compatible provider wrapper;
  \item \mir{} schema v0 and provider capability registry;
  \item adapters for OpenAI-style, Anthropic-style, and LiteLLM-compatible APIs;
  \item context block ledger, token estimation, and cost estimation;
  \item exact duplicate context elimination and stable prefix detection;
  \item prompt-cache layout reporting;
  \item simple retry, fallback, and static routing policy;
  \item execution evidence bundle and replayable trace artifact;
  \item CLI report for baseline-versus-optimized execution.
\end{itemize}

It should not lead with general semantic compression, general semantic caching, autonomous global routing, cross-provider equivalence, or claims of identical answers. These can be experimental passes later.

A credible v0 demo is a multi-step RAG or code-review workflow in which \migaki{} reduces context duplication, improves cache layout, routes cheap classification/ranking calls to cheaper models, keeps synthesis on stronger models, validates structured output, retries only failed branches, and emits an evidence report showing cost, latency, and validator deltas.
